\subsubsection*{Antes de la sesión}
\begin{itemize}
\item[$\square$] Confirmar modalidad (Teams / híbrido / presencial) y materiales.
\item[$\square$] Revisar entregables de la sesión anterior (muestra).
\item[$\square$] Preparar demo en vivo con ejemplo STEM neutro EIC.
\item[$\square$] Recordar política de integridad / uso de IA (\ceddie).
\item[$\square$] Recordar presupuesto de trabajo independiente (bloques A--D; total 16\,h).
\end{itemize}

\subsubsection*{Mensajes clave}
\begin{enumerate}
\item El outline es un \textbf{contrato pedagógico}, no un índice decorativo.
\item Menos nodos bien alineados $>$ muchos nodos genéricos.
\item El mapa de notación ahorra dolor en la sesión 4.
\end{enumerate}

\subsubsection*{Ruta sugerida (3\,h)}
\paragraph{Recap} Mapear 1 brief del grupo (con permiso) a posibles capítulos.
\paragraph{Arquitectura} Plantilla mental: Objetivo $\rightarrow$ Públicos $\rightarrow$ Capítulos $\rightarrow$ Secciones $\rightarrow$ Evaluaciones $\rightarrow$ Glosario.
\paragraph{Demo} Generar outline con restricciones de nivel; rechazar nodos vagos.
\paragraph{Taller} Priorizar 6--10 nodos; RA explícitos; mapa de símbolos.
\paragraph{Cierre} Encargar avance hacia primera sección (TI bloque C).

\subsubsection*{Contingencias}
\begin{itemize}
\item Si la herramienta falla: continuar con diseño en documento colaborativo.
\item Grupo heterogéneo: pista rápida y pista avanzada (multi-agente).
\item Dudas legales/políticas: anotar y derivar a CEDDIE; no improvisar normativa.
\end{itemize}

\subsubsection*{Checklist post-sesión}
\begin{itemize}
\item[$\square$] Recordar canal de entrega del producto de la sesión.
\item[$\square$] Asignar / recordar TI del intervalo (ver cronograma parte B).
\item[$\square$] Anotar fricciones de la demo para la siguiente sesión.
\end{itemize}
